\chapter{Introdução}
\label{chap:intro}
%\thispagestyle{plain}
\graphicspath{{./Cap1_introducao/Figures/}}

A quarta revolução industrial, denominada Indústria 4.0, representa um paradigma transformador na manufatura global, caracterizada pela integração sistemática de tecnologias digitais avançadas nos processos produtivos. Este movimento revolucionário fundamenta-se na conectividade entre sistemas físicos e digitais, proporcionando maior eficiência, flexibilidade e sustentabilidade aos processos industriais \cite{Schwab2016}. Neste contexto, a automação industrial emerge como elemento central, demandando investimentos significativos em tecnologias que possibilitem a otimização de processos e a redução de custos operacionais.

Os investimentos em tecnologia para automatização têm experimentado crescimento exponencial nas últimas décadas, impulsionados pela necessidade de competitividade global e pela busca por processos mais eficientes e sustentáveis. Segundo dados da Associação Brasileira de Automação Industrial (ABAI), o mercado brasileiro de automação industrial registrou crescimento consistente, evidenciando a importância estratégica deste setor para o desenvolvimento econômico nacional. Tais investimentos abrangem desde a modernização de equipamentos até a implementação de sistemas de controle avançados, visando à otimização de recursos e ao aumento da produtividade.

A utilização de inteligência artificial na indústria constitui uma das principais tendências da Indústria 4.0, oferecendo soluções inovadoras para desafios complexos de controle e otimização. Técnicas como redes neurais, algoritmos genéticos e lógica Fuzzy têm demonstrado eficácia significativa na resolução de problemas industriais que envolvem incertezas, não linearidades e múltiplas variáveis. A lógica Fuzzy, em particular, destaca-se pela sua capacidade de modelar conhecimento humano especializado e tratar informações imprecisas, características comuns em ambientes industriais reais.

O controle automático representa um dos pilares fundamentais da automação industrial moderna, possibilitando a manutenção de variáveis de processo dentro de faixas operacionais desejadas sem intervenção humana direta. A implementação de sistemas de controle automático resulta em benefícios tangíveis, incluindo maior precisão no controle de processos, redução de variabilidade na qualidade do produto, diminuição de custos operacionais e melhoria nas condições de segurança operacional. Ademais, tais sistemas contribuem para a redução do impacto ambiental através da otimização do uso de recursos naturais e energia.

No contexto específico da automação industrial, o controle de nível assume importância crítica em diversas aplicações, desde o armazenamento de matérias-primas até o processamento final de produtos. Sistemas de controle de nível são essenciais para garantir a continuidade operacional, prevenir perdas de produto por transbordamento, evitar danos a equipamentos por operação em vazio e assegurar a qualidade do produto final. Nas indústrias petroquímica, farmacêutica e de alimentos e bebidas, o controle preciso de nível em tanques e reatores é fundamental para manter as condições ideais de processo e atender aos rigorosos padrões de qualidade e segurança exigidos por essas indústrias.

Diante deste cenário, o presente trabalho propõe uma abordagem para o controle de nível industrial, integrando a lógica Fuzzy com o protocolo OPC, visando demonstrar a viabilidade da aplicação de técnicas de inteligência artificial em ambientes industriais reais através de uma arquitetura de integração padronizada e amplamente adotada pela indústria.

\section{Revisão Bibliográfica}
\label{sec:revbibli}

\subsection{Importância do controle de nível}
\label{sec:importancenivel}
O controle de nível em processos industriais constitui uma das variáveis mais críticas para o funcionamento seguro e eficiente de plantas produtivas. A relevância deste tipo de controle é evidenciada em diversos setores industriais, desde aplicações didáticas até implementações em larga escala industrial.

\cite{bacovis2016} realizaram um estudo comparativo entre controladores Fuzzy e PID aplicados em uma planta didática de nível de líquido, demonstrando que o controle adequado de nível é fundamental para o aprendizado e compreensão dos princípios de automação industrial. O trabalho evidencia que plantas de nível são amplamente utilizadas em ambientes educacionais devido à sua capacidade de representar fielmente os desafios encontrados em processos industriais reais, incluindo não linearidades, tempo morto e distúrbios externos.

Na indústria de alimentos e bebidas, \cite{gomes2022} investigou a aplicação da lógica Fuzzy no controle de qualidade na produção de cerveja, onde o controle preciso de nível em tanques de fermentação e maturação é crucial para garantir a qualidade do produto final. O autor demonstra que variações não controladas no nível podem afetar diretamente características organolépticas da cerveja, como aroma, sabor e teor alcoólico, evidenciando a importância crítica do controle de nível em processos biotecnológicos.

Em aplicações de saneamento e abastecimento público, \cite{silveira2021} apresentaram um estudo sobre a aplicação da lógica Fuzzy no controle de nível de reservatório de abastecimento de água, destacando que o controle eficiente de nível é essencial para garantir o fornecimento contínuo e adequado de água potável. Os autores enfatizam que falhas em sistemas de controle de nível podem resultar em desabastecimento populacional, desperdício de recursos hídricos e comprometimento da qualidade da água distribuída.

A importância do controle de nível transcende aspectos puramente técnicos, envolvendo questões econômicas, ambientais e de segurança. Sistemas de controle de nível inadequados podem resultar em perdas significativas de produto, consumo excessivo de energia, riscos ambientais por vazamentos e comprometimento da segurança operacional. Ademais, o controle preciso de nível é fundamental para otimizar o uso de recursos naturais e contribuir para a sustentabilidade de processos industriais.

\subsection{Utilização do OPC em plantas industriais}
\label{sec:utilizacaoopc}
O protocolo OPC (Open Platform Communications) e sua evolução para OPC UA (Unified Architecture) têm revolucionado a comunicação industrial, oferecendo soluções padronizadas para a integração de sistemas heterogêneos em ambientes produtivos. A aplicação desta tecnologia em plantas industriais demonstra sua versatilidade e eficácia na resolução de desafios complexos de interoperabilidade e digitalização.

\cite{carvalho2023} investigaram a digitalização de uma planta industrial utilizando o protocolo de comunicação OPC-UA, focando no sistema educacional CP Lab da Festo. O estudo demonstra como a Indústria 4.0 está impulsionando transformações profundas na indústria através de soluções que aprimoram a produção e elevam a qualidade dos produtos. Os autores evidenciam que tecnologias como IoT, sistemas ciberfísicos, Big Data e computação em nuvem estão remodelando as operações empresariais através da integração de sistemas, criação de fábricas inteligentes e melhoria da eficiência. A pesquisa destaca que a utilização do protocolo OPC UA para conectar módulos ciberfísicos e dispositivos à nuvem possibilita análise de dados em tempo real para identificar status atual, problemas potenciais e tomar decisões mais eficazes.

Em aplicações de manufatura flexível, \cite{silva2023} desenvolveu um sistema de instrumentação para uma planta de manufatura flexível utilizando o padrão OPC-UA embarcado no microcontrolador ESP32. O trabalho evidencia que o OPC-UA é um padrão de comunicação industrial que permite transferência de dados segura e confiável entre diferentes sistemas e dispositivos em setores industriais. O estudo demonstra a viabilidade de implementar servidores OPC-UA em dispositivos embarcados de baixo custo, estabelecendo comunicação entre servidores OPC-UA utilizando a ferramenta Node-RED para transferir dados do sistema de medição para controladores lógicos programáveis (CLP) presentes na planta industrial.

\cite{petrocchi2024} apresentou uma implementação da integração entre tecnologia RFID e o padrão OPC UA aplicada a um sistema de manufatura flexível no laboratório da UNESP campus Sorocaba. O trabalho demonstra que o padrão OPC UA fornece uma troca segura de informações e dados na área da automação industrial, utilizando um protocolo padrão extensível que pode ser empregado em multiplataformas. A pesquisa evidencia que é possível realizar alterações das informações através do sistema de servidor e cliente graças às funcionalidades de acesso aos dados do OPC UA. O desenvolvimento utilizou o ambiente Node-RED para configuração automática do leitor RFID e criação do servidor OPC UA, permitindo leitura e escrita de dados tanto pelo servidor quanto pela interface gráfica.

\cite{souza2024} conduziu um estudo de caso sobre controle e supervisão de processos utilizando o padrão OPC, focando na superação dos desafios de integração entre equipamentos e softwares de diferentes fabricantes devido a protocolos proprietários. O trabalho utilizou o CLP Codesys e o ambiente simulado do Factory IO, com foco na planta industrial Sorting Station, demonstrando como o padrão OPC UA promove a troca de dados e comandos entre sistemas distintos. A pesquisa explora as possibilidades tecnológicas oferecidas por essa integração, evidenciando a capacidade do OPC UA em conectar e comunicar softwares e dispositivos de diferentes fabricantes.

A utilização do OPC em plantas industriais transcende aspectos puramente técnicos, representando um facilitador fundamental para a implementação da Indústria 4.0. Os trabalhos analisados demonstram que a adoção do protocolo OPC UA resulta em maior flexibilidade operacional, redução de custos de integração, melhoria na interoperabilidade entre sistemas e facilita a migração para arquiteturas mais modernas. Ademais, a capacidade do OPC UA de operar em dispositivos embarcados de baixo custo democratiza o acesso a tecnologias avançadas de comunicação industrial, possibilitando a modernização de plantas existentes com investimentos reduzidos.

\subsection{Controle inteligente e OPC em aplicações industriais}
\label{sec:controleintopc}
A convergência entre técnicas de controle inteligente e protocolos de comunicação padronizados como o OPC representa uma das principais tendências na automação industrial moderna. Esta integração possibilita o desenvolvimento de sistemas de controle mais sofisticados, flexíveis e adaptáveis às demandas crescentes da Indústria 4.0, oferecendo soluções inovadoras para desafios complexos de controle e manutenção industrial.

\cite{oliveira2023} desenvolveu um estudo sobre gerenciamento de nível em reservatório de líquidos utilizando lógica Fuzzy e controle PID, implementando uma arquitetura integrada baseada em OPC para comunicação entre sistemas. O trabalho apresenta um sistema amplamente utilizado nas instalações industriais: o controle de nível de tanque, realizado através do software CodeSys e Matlab/Simulink, com simulação no ambiente Factory I/O. A pesquisa desenvolveu dois controladores distintos: um baseado em lógica Fuzzy e outro em controle PID, utilizando o protocolo de comunicação OPC para estabelecer a comunicação entre o CLP, Matlab e Factory I/O. O autor evidencia que a utilização da ferramenta Fuzzy Logic Designer do Matlab, em conjunto com o Simulink, proporcionou análise detalhada do comportamento do sistema, demonstrando controle satisfatório do sistema simulado de tanque. Esta abordagem integrada ilustra como a combinação de técnicas de controle inteligente com protocolos de comunicação padronizados pode resultar em sistemas de controle mais eficazes e facilmente integráveis a arquiteturas industriais existentes.

\cite{coretti2025} investigou a automação com redes inteligentes para manutenção de sistemas de controle de processos industriais, focando no desenvolvimento de mecanismos para digitalização da malha de controle de vazão e tomada de decisão para análise da manutenção de sistemas automatizados. O trabalho utiliza conceitos da Indústria 4.0 e algoritmos de inteligência artificial para melhorar a visibilidade das malhas de controle fechadas e dos ativos da fábrica. A pesquisa propõe uma abordagem que facilita à equipe industrial de manutenção e operação maior facilidade na gestão dos sistemas de automação e tomada de decisões precisas, sem grandes interrupções do processo produtivo. O autor demonstra que a implementação de redes industriais inteligentes, combinada com manutenção proativa e softwares industriais, permite gestão mais eficiente das fábricas com custos reduzidos. A proposta inclui interface de digitalização da malha de controle livre, configurada em blocos funcionais na web, evidenciando a tendência de integração entre sistemas de automação tradicionais e plataformas digitais modernas.

A integração entre controle inteligente e OPC em aplicações industriais representa uma evolução natural dos sistemas de automação, proporcionando benefícios significativos em termos de flexibilidade, manutenibilidade e capacidade de integração. Os trabalhos analisados demonstram que esta abordagem permite não apenas melhor desempenho de controle através de técnicas inteligentes como lógica Fuzzy, mas também facilita a implementação de estratégias de manutenção preditiva e proativa através da padronização de comunicação proporcionada pelo OPC.

Ademais, a utilização de protocolos OPC em sistemas de controle inteligente contribui para a democratização do acesso a tecnologias avançadas, permitindo que empresas de diferentes portes implementem soluções sofisticadas de automação sem dependência de fornecedores específicos. Esta característica é fundamental para a competitividade industrial no contexto da Indústria 4.0, onde a capacidade de adaptação rápida às mudanças tecnológicas e de mercado constitui vantagem estratégica crítica.

\section{Objetivos}
\label{sec:obj}
Os objetivos deste trabalho podem ser definidos da seguinte forma:

\subsection{Objetivo Geral}
\label{sec:objgeral}
Desenvolver e comparar o desempenho de dois métodos de controle, PID e PID-Fuzzy, em uma planta de controle de luminosidade, a fim de identificar qual método oferece melhor resposta em termos de estabilidade, tempo de resposta e precisão.

\subsection{Objetivos Específicos}
\label{sec:objespec}
Para alcançar o objetivo geral proposto, foram definidos os seguintes objetivos específicos:

\begin{itemize}
\item Caracterização do sistema de controle: Extrair os parâmetros Fuzzy da planta de nível, identificando e definindo as variáveis de entrada e saída do sistema, bem como suas respectivas funções de pertinência e regras de inferência necessárias para o controle eficaz do processo.
\item Desenvolvimento da API de controle: Desenvolver uma API especializada utilizando a linguagem Python, implementando os algoritmos de lógica Fuzzy baseados nos parâmetros previamente definidos, garantindo flexibilidade, robustez e facilidade de integração com sistemas externos.
\item Implementação da aplicação de integração: Desenvolver uma aplicação utilizando a linguagem C-Sharp para promover a integração entre a API de controle Fuzzy e a planta de nível, implementando os protocolos de comunicação OPC necessários para estabelecer a troca de dados em tempo real entre os sistemas.
\item Desenvolvimento da interface humano-máquina: Criar uma IHM (Interface Humano-Máquina) intuitiva e funcional que permita a configuração dinâmica da API e do cliente OPC, bem como a visualização em tempo real do controle de nível, proporcionando monitoramento eficiente e facilidade de operação do sistema integrado.
\item Análise de desempenho: Analisar os resultados obtidos através de testes e simulações, avaliando parâmetros como estabilidade do controle, tempo de resposta, precisão no seguimento de referência, robustez a distúrbios e eficácia da comunicação OPC, documentando as conclusões e contribuições do trabalho para a área de automação industrial.
\end{itemize}

\section{Estrutura do Texto}
Este trabalho está organizado em seis capítulos, estruturados de forma a apresentar uma sequência lógica e progressiva do desenvolvimento da pesquisa, desde a fundamentação teórica até a análise dos resultados obtidos.

\begin{itemize}
\item Capítulo 1 - Introdução: Apresenta o contexto geral da pesquisa, abordando a importância do controle de nível em plantas industriais, a utilização do protocolo OPC em ambientes produtivos e a convergência entre controle inteligente e comunicação OPC. Este capítulo também define os objetivos do trabalho e estabelece a justificativa para a integração de técnicas de lógica Fuzzy com protocolos de comunicação padronizados.
\item Capítulo 2 - Fundamentação Teórica: Explora a evolução histórica dos sistemas de controle, desde as técnicas clássicas até as abordagens modernas, destacando onde o controle Fuzzy se posiciona nesta evolução tecnológica. Adicionalmente, examina a utilização de protocolos de comunicação em plantas industriais, enfatizando a importância da padronização e interoperabilidade para a Indústria 4.0.
\item Capítulo 3 - Metodologia: Detalha a metodologia empregada no desenvolvimento do trabalho, incluindo a justificativa para a escolha da planta de nível como objeto de estudo, a seleção da tecnologia Python para criação da API de controle Fuzzy, e a definição da tecnologia para desenvolvimento da Interface Humano-Máquina (IHM), fundamentando cada decisão tecnológica com base em critérios técnicos e práticos.
\item Capítulo 4 - Desenvolvimento: Apresenta o processo de desenvolvimento dos componentes principais do sistema integrado, incluindo a implementação da API de controle Fuzzy, o desenvolvimento da IHM para monitoramento e configuração, e a criação do cliente OPC UA para estabelecer a comunicação entre os diferentes módulos do sistema.
\item Capítulo 5 - Resultados e Discussão: Analisa os resultados obtidos através da integração entre a API de controle Fuzzy e a planta de nível via protocolo OPC, apresentando métricas de desempenho, estabilidade do sistema e eficácia da comunicação, bem como a discussão crítica dos resultados em relação aos objetivos propostos.
\item Capítulo 6 - Conclusões e Trabalhos Futuros: Sintetiza as principais conclusões do trabalho, destacando as contribuições para a área de automação industrial e controle inteligente, além de propor direcionamentos para trabalhos futuros que possam ampliar e aprofundar os resultados obtidos nesta pesquisa.
\end{itemize}